\documentclass{beamer}

\mode<presentation>





\usetheme{Frankfurt}%
\usecolortheme{seagull}
\logo{\includegraphics[height=.25in]{clarksonGreen}}

\definecolor{garnet}{RGB}{136,0,0}
%\definecolor{clarksonGreen}{RGB}{0,71,28}
\definecolor{clarksonGreen}{RGB}{0,52,21}
\setbeamercolor{palette primary}{fg=clarksonGreen,bg=white}
\setbeamercolor{palette secondary}{fg=clarksonGreen,bg=white}
\setbeamercolor{palette tertiary}{fg=clarksonGreen,bg=white}
\setbeamercolor{palette quaternary}{bg=clarksonGreen,fg=white}
\setbeamercolor{block title}{fg=black,bg=black!15}
\setbeamercolor{block body}{fg=black,bg=black!10}
\setbeamercolor{titlelike}{bg=clarksonGreen,fg=white} % parent=palette quaternary}

\newcommand{\half}{\mbox{$\frac{1}{2}$}}
\newcommand{\deltat}{\mbox{$\triangle t$}}
\newcommand{\deltax}{\mbox{$\triangle x$}}
\newcommand{\deltay}{\mbox{$\triangle y$}}

\newcommand{\deriv}[2]{\frac{d}{d#2}#1}
\newcommand{\derivTwo}[2]{\frac{d^2}{d#2^2}#1}

\newcommand{\lp}{\left(}
\newcommand{\rp}{\right)}



\begin{document}

\title{Ordinary Differential Equations}
\subtitle{Math 232 - Week 8, Day 2}

\author{Kelly Black}
\institute{Clarkson University}
\date{19 Oct 2011}

\begin{frame}
  \titlepage
\end{frame}

\begin{frame}
  \frametitle{Outline}
  \tableofcontents[pausesection,hideallsubsections]
\end{frame}


\section{Non-homogeneous, Constant Coefficient, Second Order Equations}


\begin{frame}
  \frametitle{Non-homogeneous, Constant Coefficient, Second Order
    Equations}
  What if the equation is \textbf{not} homogeneous?

  Example: An horizontal spring/mass system that has a constant force
  of 1N to the right.

  \begin{eqnarray*}
    \frac{d}{dt} \lp m v_x \rp & = & -b v_x - kx + 1, \\
    m x'' + bx' + kx & = & 1.
  \end{eqnarray*}

  What do we do?


\end{frame}


\begin{frame}
  \frametitle{Note About Linear Equations}

  These are linear equations. If $y_h(t)$ is a solution to the
  homogeneous problem and $y_p(t)$ is a solution to the particular
  problem then
  \begin{eqnarray*}
    y_h(t) + y_p(t)
  \end{eqnarray*}
  is also a solution.

\end{frame}


\begin{frame}
  \frametitle{Overview of the Solution Method}

  Given
  \begin{eqnarray*}
    a y'' + by' + cy & = & f(t):
  \end{eqnarray*}
  \begin{enumerate}
  \item Solve the homogeneous problem,
    \begin{eqnarray*}
      a y_h'' + by_h' + cy_h & = & 0.
    \end{eqnarray*}
  \item Solve the particular solution,
    \begin{eqnarray*}
      a y_p'' + by_p' + cy_p & = & f(t).
    \end{eqnarray*}
  \item Form the solution
    \begin{eqnarray*}
      y(t) & = & y_h(t) + y_p(t).
    \end{eqnarray*}
  \end{enumerate}


\end{frame}

\section{Examples}

\begin{frame}
  \frametitle{Example}

  \begin{eqnarray*}
    y'' + 4y' + 4y & = & 1.
  \end{eqnarray*}

  \uncover<2->
  {
    Find the homogeneous solution:
      \begin{eqnarray*}
        r^2 + 4r + 4 & = &, \\
        (r+2)^2 & = & 0, \\
        y_h & = & C_1 e^{-2t} + C_2 t e^{-2t}.
      \end{eqnarray*}
    }

\end{frame}


\begin{frame}
  \frametitle{The Particular Solution}
  
  We make a guess:
  \begin{eqnarray*}
    y_p & = & A, \\
    y_p' & = 0, \\
    y_p'' & = & 0, \\
    0 + 4(0) + 4A & = & 1, \\
    A & = & \frac{1}{4}, \\
    y_p & = & \frac{1}{4}
  \end{eqnarray*}

  The solution to the problem is
  \begin{eqnarray*}
    y(t) & = & C_1 e^{-2t} + C_2 t e^{-2t} + \frac{1}{4}.
  \end{eqnarray*}

\end{frame}

\section{General Idea}

\begin{frame}
  \frametitle{General Idea}

  \begin{enumerate}
  \item If $y(t)$ is a polynomial of degree $n$ then
    \begin{eqnarray*}
      a y'' + b y' + cy 
    \end{eqnarray*}
    is also a polynomial of degree $n$.
  \item If $y(t)$ is an exponential function then
    \begin{eqnarray*}
      a y'' + by' + cy 
    \end{eqnarray*}
    is also an exponential function.
  \item If $y(t)$ is made up of sines \textbf{or} cosines then
    \begin{eqnarray*}
      a y'' + by' + cy
    \end{eqnarray*}
    is  made up of sines \textbf{and} cosines.
  \end{enumerate}

\end{frame}


\begin{frame}
  \frametitle{Example}

  \begin{eqnarray*}
    y'' + 4y' + 4y & = & 2t^2.
  \end{eqnarray*}

  \uncover<2->
  {
    \begin{eqnarray*}
      y_h & = & C_1 e^{-2t} + C_2 t e^{-2t}
    \end{eqnarray*}
  }

\end{frame}


\begin{frame}
  \frametitle{The particular solution}

  \begin{eqnarray*}
    y_p & = & at^2 + bt + c, \\
    y_p' & = & 2at + b \\
    y_p'' & = & 2a
  \end{eqnarray*}

  Substitute this back into the equation:
  \begin{eqnarray*}
    2a + 4 (2at+b) + 4 (at^2 + bt + c) & = & 2t^2, \\
    \Rightarrow 
    a & = & \half \\
    b & = & -1 \\
    c & = & \frac{3}{4}
  \end{eqnarray*}

  The solution is 
  \begin{eqnarray*}
    y(t) & = & C_1 e^{-2t} + C_2 t e^{-2t} + \half t^2 - t + \frac{3}{4}.
  \end{eqnarray*}

\end{frame}


\begin{frame}
  \frametitle{Example}

  \begin{eqnarray*}
    y'' + 4y' + 4y & = & 2 e^{3t}.
  \end{eqnarray*}

  \uncover<2->
  {
    \begin{eqnarray*}
      y_h & = & C_1 e^{-2t} + C_2 t e^{-2t}
    \end{eqnarray*}
  }

\end{frame}


\begin{frame}
  \frametitle{The particular solution}

  \begin{eqnarray*}
    y_p & = & A e^{3t}, \\
    y_p' & = & 3 A e^{3t} \\
    y_p'' & = & 9 A e^{3t}
  \end{eqnarray*}

  Substitute this back into the equation:
  \begin{eqnarray*}
    9 A e^{3t} + 4 (3 A e^{3t}) + 4 (A e^{3t}) & = & 2 e^{3t}, \\
    \Rightarrow 
    A & = & \frac{2}{25}
  \end{eqnarray*}

  The solution is 
  \begin{eqnarray*}
    y(t) & = & C_1 e^{-2t} + C_2 t e^{-2t} + \frac{2}{25} e^{3t}.
  \end{eqnarray*}

\end{frame}



\begin{frame}
  \frametitle{Example}

  \begin{eqnarray*}
    y'' + 4y' + 4y & = & 2 \cos(3t).
  \end{eqnarray*}

  \uncover<2->
  {
    \begin{eqnarray*}
      y_h & = & C_1 e^{-2t} + C_2 t e^{-2t}
    \end{eqnarray*}
  }

\end{frame}


\begin{frame}
  \frametitle{The particular solution}

  \begin{eqnarray*}
    y_p & = & a \cos(3t) + b \sin(3t), \\
    y_p' & = & -3 a \sin(3t) + 3b \cos(3t) \\
    y_p'' & = & -9 a \cos(3t) - 9b \sin(3t)
  \end{eqnarray*}

  Substitute this back into the equation:
  \begin{eqnarray*}
    -9 a \cos(3t) - 9b \sin(3t) 
    + 4 \lp -3 a \sin(3t) + 3b \cos(3t) \rp & & \\
    + 4 \lp a \cos(3t) + b \sin(3t) \rp
     & = & 2 \cos(3t), \\
    \Rightarrow 
    a & = & \frac{-10}{169} \\
    b & = & \frac{24}{169}
  \end{eqnarray*}

  The solution is 
  \begin{eqnarray*}
    y(t) & = & C_1 e^{-2t} + C_2 t e^{-2t} + \frac{-10}{169}\cos(3t) +
    \frac{24}{169} \sin(3t).
  \end{eqnarray*}

\end{frame}


\begin{frame}
  \frametitle{Example}

  \begin{eqnarray*}
    y'' + 4y' + 4y & = & 2t^2 + 2 e^{3t}.
  \end{eqnarray*}

  \uncover<2->
  {
    \begin{eqnarray*}
      y_h & = & C_1 e^{-2t} + C_2 t e^{-2t}
    \end{eqnarray*}
  }

\end{frame}


\begin{frame}
  \frametitle{The particular solution}

  \begin{eqnarray*}
    y_p & = & at^2 + bt + c + d e^{3t}, \\
    y_p' & = & 2at + b + 3e^{3t} \\
    y_p'' & = & 2a + 9 e^{3t}
  \end{eqnarray*}

  Substitute this back into the equation:
  \begin{eqnarray*}
    2a + 9de^{3t} + 4 (2at+b + 3 e^{3t}) + 4 (at^2 + bt + c + e^{3t}) & = & 2t^2+2e^{3t}, \\
    \Rightarrow 
    a & = & \half \\
    b & = & -1 \\
    c & = & \frac{3}{4} \\
    d & = & \frac{2}{25} e^{3t}
  \end{eqnarray*}

  The solution is 
  \begin{eqnarray*}
    y(t) & = & C_1 e^{-2t} + C_2 t e^{-2t} + \half t^2 - t +
    \frac{3}{4} + \frac{2}{25} e^{3t}R.
  \end{eqnarray*}

\end{frame}



\begin{frame}
  \frametitle{Example}

  \begin{eqnarray*}
    y'' + 4y' + 4y & = & 2 \cos(3t) + 2e^{3t}.
  \end{eqnarray*}

  \uncover<2->
  {
    \begin{eqnarray*}
      y_h & = & C_1 e^{-2t} + C_2 t e^{-2t}
    \end{eqnarray*}
  }

\end{frame}


\begin{frame}
  \frametitle{The particular solution}

  \begin{eqnarray*}
    y_p & = & a \cos(3t) + b \sin(3t) + ce^{3t}, \\
    y_p' & = & -3 a \sin(3t) + 3b \cos(3t) + 3 c e^{3t}\\
    y_p'' & = & -9 a \cos(3t) - 9b \sin(3t) + 9 c e^{3t}
  \end{eqnarray*}

  Substitute this back into the equation:
  \begin{eqnarray*}
    -9 a \cos(3t) - 9b \sin(3t) + 9 c e^{3t} & & \\
    + 4 \lp -3 a \sin(3t) + 3b \cos(3t) + 3c e^{3t}\rp  & & \\
    + 4 \lp a \cos(3t) + b \sin(3t) + c e^{3t}\rp
     & = & 2 \cos(3t), \\
    \Rightarrow 
    a & = & \frac{-10}{169} \\
    b & = & \frac{24}{169} \\
    c & = & \frac{2}{25} 
  \end{eqnarray*}


\end{frame}


\begin{frame}

  The solution is 
  \begin{eqnarray*}
    y(t) & = & C_1 e^{-2t} + C_2 t e^{-2t} + \frac{-10}{169}\cos(3t) +
    \frac{24}{169} \sin(3t) + \frac{2}{25}  e^{3t}.
  \end{eqnarray*}
  
\end{frame}



\end{document}

% LocalWords:  Clarkson pausesection hideallsubsections
